\documentclass[12pt]{article}

\usepackage{amsmath,amssymb,amsthm}
\usepackage{geometry}
\usepackage{enumitem}
\usepackage{longtable}
\usepackage{booktabs}
\usepackage{hyperref}
\usepackage{array}
\usepackage{microtype}
\geometry{margin=1in}
\hypersetup{pdfborder={0 0 0}}
\title{Unus Solus Possibilis Est\\
\large Corpus Closure \emph{and} Uniqueness of the Admissible Interior\\}
\author{Amos Jay Maley}
\date{February 9, 2026}

\newtheorem{definition}{Definition}
\newtheorem{axiom}{Axiom}
\newtheorem{lemma}{Lemma}
\newtheorem{proposition}{Proposition}
\newtheorem{theorem}{Theorem}
\newtheorem{corollary}{Corollary}
\newtheorem{remark}{Remark}

\begin{document}
\maketitle

\begin{abstract}
This paper is a standalone derivation of the apex closure claim of the Maley corpus at the apex consequence layer.
Working conditionally within a standing-instantiated application regime, and from a minimal internal axiom set---intelligibility, non-degenerate irreversible commitment, unrestricted composition,
scope discipline, meta-invariance of admissibility criteria, a \emph{binary and non-eventful} AMetric boundary,
and bivalence without metricity---we derive an internally closed standing regime and prove uniqueness of the
admissible interior (including explicit elimination of \emph{discrete} multiplicity).
Every entry in the degeneracy-block matrix is justified by in-paper lemma/theorem pointers; Appendix~A supplies
short proofs for each degeneracy class, and Appendix~B supplies an axiom minimality/independence analysis (dependency matrix,
independence witnesses, and controlled weakenings).
\end{abstract}

\tableofcontents

\section{Scope, Typing, and Fail-Closed Defaults}
\label{sec:scope}

This paper is \emph{fail-closed}. If the typed components below are not instantiated,
standing defaults to \emph{undefined}. ``Undefined'' is not a value in $[0,1]$; it is
\emph{non-instantiation} of the eligibility relation.

\begin{definition}[Acts]
An \emph{act} $a$ is any unit that purports to be standing-bearing: asserting a claim as truth-apt,
invoking an admissibility criterion, transporting confirmation, upgrading performance to authority,
transferring standing across domains, or proposing a regime translation.
\end{definition}

\begin{definition}[Typed domains]
A \emph{domain} $D$ is a typed envelope
\[
D := (E,\ A,\ F,\ C,\ \mathrm{DAC}),
\]
where:
\begin{itemize}[leftmargin=2.2em]
  \item $E$ is the domain's admissible evidence carriers (reports, records, reference-bearers),
  \item $A$ is the admissibility criterion class governing allowed transformations and descriptions,
  \item $F$ is the domain's reference-fixation structure (what reports are about, and how this is preserved),
  \item $C$ is the domain's truth-apt discriminator (confirmation discipline, if any),
  \item $\mathrm{DAC}$ is domain-of-application closure (derived, not primitive).
\end{itemize}
\end{definition}

\begin{definition}[Derived DAC (not primitive)]
$\mathrm{DAC}$ holds in $D=(E,A,F,C,\mathrm{DAC})$ iff:
\begin{enumerate}[label=(\alph*), leftmargin=2.2em]
  \item (\emph{Evidence closure}) admissible redescriptions/transport map admissible reports to admissible reports (the admissible orbit of any $e\in E$ remains in $E$);
  \item (\emph{Fixation totality}) $F$ is well-typed and invariant on the admissible orbit (no retargeting required to maintain aboutness);
  \item (\emph{Discriminator totality}) $C$ is defined on the admissible orbit without importing post-hoc restrictions.
\end{enumerate}
Thus $\mathrm{DAC}$ abbreviates ``$(E,A,F,C)$ closes on itself as a standing-bearing envelope.''
\end{definition}

\begin{definition}[Domain individuation (non-stipulative)]
Two domains $D_1=(E_1,A_1,F_1,C_1,\mathrm{DAC}_1)$ and $D_2=(E_2,A_2,F_2,C_2,\mathrm{DAC}_2)$ are distinct iff
at least one component differs in a way that changes well-formedness, reference eligibility, or truth-apt discrimination.
Topic labels, vocabulary shifts, and application rhetoric do not individuate domains.
\end{definition}

\begin{remark}[Application-regime anchor]
\label{rem:app-anchor}
The present paper is standalone only at the apex consequence layer. It does not attempt a lower derivation of admissibility, standing, reference, or irreversibility.
Its application regime is the class of typed, standing-governed regimes in which non-degenerate reasoning is already operative.
In the companion manuscripts \emph{Foundational Closure of Admissibility and Standing} (especially Theorem~2.14) and \emph{Necessity of Admissibility in Non-Degenerate Compositional Systems}
(especially Theorem~5.7, Corollary~5.8, and Theorem~7.5), kernel instantiation, fixed-domain gate existence, and standing-instantiation are established at the upstream layers.
The present paper therefore takes the standing-instantiated admissible side as fixed at application and derives the apex closure consequences internal to that already-instantiated regime.
All existence claims in this manuscript are conditional on operation within that meaningful, non-degenerate application class; they are not offered here as a fresh lower derivation independent of the upstream kernel-instantiation results.
\end{remark}

\section{Axioms (Minimal Internal Premises of Application)}
\label{sec:axioms}

The apex theorem is conditional on the following internal axioms once the standing-instantiated application regime of Remark~\ref{rem:app-anchor} is fixed.

\begin{axiom}[Intelligibility]
\label{ax:intell}
Meaningful discourse is intended: claim identity and evidential reference are required to be well-defined under admissible operations.
\end{axiom}

\begin{axiom}[Non-degenerate irreversible commitment]
\label{ax:irr}
The system contains irreversible commitment operations and intends non-degenerate reasoning under such operations. In particular,
illicit acts cannot be repaired post hoc \emph{as the same act}.
\end{axiom}

\begin{axiom}[Unrestricted composition]
\label{ax:comp}
Admissible operations compose; ad hoc restrictions may not be introduced solely to block laundering.
\end{axiom}

\begin{axiom}[Scope discipline]
\label{ax:scope}
No post-hoc narrowing/expansion/reframing is permitted to legalize otherwise-illicit moves. Any envelope shift used as repair is inadmissible.
\end{axiom}

\begin{axiom}[Meta-invariance of admissibility criteria]
\label{ax:miax}
Admissibility criteria must be invariant under admissible redescription. Any criterion depending on privileged representations or silent redescription is null.
\end{axiom}

\begin{axiom}[Binary, non-eventful AMetric boundary]
\label{ax:ametric}
The boundary separating pre-admissible proposal space from admissible reasoning is \emph{binary} and \emph{non-eventful}.
Binary: exactly two statuses exist (admissible / not admissible); no graded proximity or partial crossing.
Non-eventful: crossing is not a process or trajectory.
Therefore the boundary admits no admissibility-relevant parameterization, including discrete indexing of multiple admissible alternatives.
\end{axiom}

\begin{axiom}[Bivalence without metricity]
\label{ax:biv}
At this layer, eligibility statuses are bivalent: defined vs undefined, legitimate vs illegitimate.
There is no ``almost'' legitimacy, no proximity-to-admissible, and no graded repair regime.
\end{axiom}

\section{Standing and Gates: Core Lemma Chain}
\label{sec:gates}

\subsection{Standing as a determinate evaluation}

\begin{definition}[Standing and gate predicates]
\label{def:standing}
For each typed domain $D$, a \emph{gate predicate} is a predicate $G_D(a)$ intended to carve the acts into eligible vs ineligible acts.
When $G_D$ is instantiated, standing is the induced bivalent evaluation
\[
L_D(a)=1 \iff G_D(a), \qquad L_D(a)=0 \iff \neg G_D(a).
\]
If $G_D$ is not instantiated (or $D$ is not typed), standing in $D$ is \emph{undefined}.
\end{definition}

\begin{lemma}[Gate existence is forced by irreversibility + intelligibility]
\label{lem:gate-existence}
Under Axioms~\ref{ax:intell} and \ref{ax:irr}, any non-degenerate domain must instantiate some gate predicate $G_D$.
\end{lemma}

\begin{proof}
Axiom~\ref{ax:irr} demands that some acts be \emph{irreversibly} excluded: once excluded, they cannot become licit again \emph{as the same act}.
For that demand to be meaningful, the system must be able to (i) \emph{state} which acts are excluded and (ii) preserve that status under admissible operations.
If no gate predicate exists, then there is no internal predicate that distinguishes licit from illicit acts; exclusion cannot be expressed as a stable status.
Then either: (a) every act is treated as licit (irreversibility is void), or (b) no act is eligible (discourse collapses), or (c) eligibility varies
silently by representation, destroying claim identity and evidential reference (violating Axiom~\ref{ax:intell}). In all cases, non-degenerate reasoning fails.
Therefore, any non-degenerate domain must instantiate some gate predicate.
\end{proof}

\begin{lemma}[Binary standing (when defined)]
\label{lem:binary-standing}
Under Axiom~\ref{ax:biv}, whenever standing in $D$ is defined, $L_D(a)\in\{0,1\}$ for all acts $a$.
\end{lemma}

\begin{proof}
Axiom~\ref{ax:biv} forbids graded legitimacy. Thus the only eligible statuses (when standing is instantiated) are ``legitimate'' and ``illegitimate.''
These correspond to the two truth values of the gate predicate. Hence $L_D(a)\in\{0,1\}$.
\end{proof}

\subsection{Binary AMetric boundary blocks parameterization}

\begin{lemma}[No admissibility-relevant parameterization (including discrete indexing)]
\label{lem:no-parameter}
Under Axiom~\ref{ax:ametric}, the admissibility boundary admits no admissibility-relevant parameterization, including discrete indexing of multiple admissible alternatives.
\end{lemma}

\begin{proof}
Suppose, for contradiction, that there exists an admissibility-relevant parameterization by an index set $I$ (possibly finite) such that ``admissible'' crossing
comes in distinct indexed forms. Then (i) there must be a boundary-authoritative differentiator $\Delta(i)$ that selects which indexed admissible form applies, and (ii) that differentiator must matter for standing/admissibility.
But then the boundary has at least three distinct boundary-relevant statuses: ``inadmissible'', ``admissible of type $i$'', and ``admissible of type $j\neq i$''.
This contradicts boundary binarity. Moreover, a selection of $i$ would constitute a boundary transition mechanism, contradicting non-eventfulness.
Therefore no such parameterization exists.
\end{proof}

\begin{remark}[Consequence]
Lemma~\ref{lem:no-parameter} blocks not only continuous ``degrees'' of admissibility but also discrete plural admissible alternatives.
This is the crucial ingredient in eliminating categorical multiplicity at the apex.
\end{remark}

\begin{lemma}[Uniqueness of admissibility classification (no competing gates)]
\label{lem:unique-gate}
Assume Axioms~\ref{ax:intell}, \ref{ax:ametric}, and \ref{ax:biv}. Let $G$ and $G'$ be two candidate gate predicates for the same typed domain $D$
that are each claimed to determine admissibility/standing for acts in $D$. If there exists an act $a$ with $G(a)\neq G'(a)$, then either
standing in $D$ is indeterminate (violating intelligibility) or the boundary admits an admissibility-relevant differentiator selecting between
$G$ and $G'$ (contradicting Lemma~\ref{lem:no-parameter}). Therefore, within the intended regime, admissibility classification for a fixed domain is
unique up to extensional equivalence of the gate predicate.
\end{lemma}

\begin{proof}
If $G(a)\neq G'(a)$ for some act $a$, then the system has two incompatible admissibility classifications for the same act in the same domain.
To maintain determinacy of standing (Axiom~\ref{ax:intell}), the system must specify which classification governs eligibility. Any such specification
is an admissibility-relevant differentiator: it selects between admissibility outcomes for $a$. By Lemma~\ref{lem:no-parameter}, admissibility-relevant
differentiators (including discrete index choices) are forbidden by the binary, non-eventful AMetric boundary. Hence no such disagreement is permissible.
Therefore any two admissibility predicates for a fixed domain must coincide extensionally.
\end{proof}


\section{Kernel Closure: No Repairs, No Generators, No Carriers}
\label{sec:kernel}

\subsection{No repairs}

\begin{theorem}[No repairs for an evaluated act]
\label{thm:no-repairs}
Assume standing is defined in $D$. If $L_D(a)=0$, then no internal post-processing operation on the \emph{same act} $a$
yields $L_D(a)=1$. Any legitimate ``repair'' must produce a new act $a'$ that independently satisfies $G_D$.
\end{theorem}

\begin{proof}
Let $a$ be an evaluated act with $L_D(a)=0$. Suppose an internal operation $R$ purports to make $a$ legitimate.
There are only two cases.

\textbf{Case 1 (identity-preserving).} $R$ preserves the identity of $a$ as an evaluated act (it is ``the same act'').
Then $R$ changes $L_D(a)$ from $0$ to $1$ for the same act, which is exactly the post-hoc reversal forbidden by Axiom~\ref{ax:irr}.
Therefore identity-preserving repair is impossible.

\textbf{Case 2 (identity-changing).} $R$ produces a distinct act $a'$. Then $a'$ must be evaluated separately.
By Definition~\ref{def:standing}, $L_D(a')=1$ holds iff $G_D(a')$ holds. Thus ``repair'' produces a new candidate act whose legitimacy is not inherited
from $a$ but depends on fresh gating. In either case, the original act $a$ remains illegitimate.
\end{proof}

\subsection{No generators}

\begin{theorem}[No generators]
\label{thm:no-generators}
There is no operation $T$ on acts such that, for some act $a$ with $L_D(a)=0$, one has $L_D(T(a))=1$ \emph{as a consequence of $a$ alone}
without $T(a)$ constituting a fresh gated act in $D$.
\end{theorem}

\begin{proof}
Let $a$ satisfy $L_D(a)=0$ and consider any operation $T$.

If $T(a)$ is the same evaluated act as $a$, then $T$ is an identity-preserving ``repair,'' ruled out by Theorem~\ref{thm:no-repairs}.
If $T(a)$ is a distinct act $a'$, then legitimacy of $a'$ is not a function of $a$ alone: $L_D(a')=1$ holds only if $G_D(a')$ holds.
Thus there is no generator that converts an illegitimate act into a legitimate act \emph{by transformation alone}.
\end{proof}

\subsection{No carriers (numeric \emph{or} structural)}

\begin{definition}[Carrier attempt]
A \emph{carrier attempt} is any derived predicate $Q_D(a)$ proposed to substitute for gating as the primitive authorizer of standing,
i.e.\ a rule of the form ``$L_D(a)=1$ because $Q_D(a)$'', where $Q_D$ is not merely a renaming of $G_D$.
Carrier attempts include: scalar rules ($q>\tau$), probabilities and typicalities, confidence scores, losses/utility, information criteria,
and also non-scalar invariants (algebraic, categorical, symmetry, embedding, etc.) proposed to authorize standing.
\end{definition}

\begin{lemma}[Any primitive authorizer is a gate predicate]
\label{lem:authorizer-is-gate}
If $Q_D$ is proposed to primitively authorize standing, then $Q_D$ is, extensionally, a gate predicate: it partitions acts into eligible vs ineligible.
\end{lemma}

\begin{proof}
Primitive authorization means: the system uses $Q_D$ to decide which acts are eligible. That decision is a partition of acts into two classes.
Under bivalence, there is no third class. Therefore $Q_D$ functions as a gate predicate, regardless of its internal form.
\end{proof}

\begin{theorem}[No carriers (numeric or structural)]
\label{thm:no-carriers}
Assume Axioms~\ref{ax:intell}, \ref{ax:miax}, \ref{ax:ametric}, and \ref{ax:biv}. Then no carrier attempt $Q_D$ can serve as a \emph{primitive} authorizer of standing (i.e.\ no $Q_D\not\equiv G_D$ can decide eligibility) except by coinciding extensionally with the unique gate predicate guaranteed by Lemma~\ref{lem:unique-gate}.
Any admissible use of $Q_D$ must be downstream bookkeeping inside an already gated envelope and must have no authority over admissibility/standing.
\end{theorem}

\begin{proof}
Let $Q_D$ be a carrier attempt. By Lemma~\ref{lem:authorizer-is-gate}, if $Q_D$ is used primitively to decide eligibility, then $Q_D$ is functioning as a gate predicate.

We show that any nontrivial carrier attempt cannot be a legitimate primitive authorizer.

\textbf{(1) Metric/probabilistic carriers.} If $Q_D$ is defined via metric proximity, measure, typicality, probabilities, rankings, convergence, or any
``almost'' predicate that is meant to \emph{decide} eligibility, then admissibility becomes transmissive in metric/probabilistic structure.
This contradicts Axiom~\ref{ax:ametric} (binary, non-eventful AMetric boundary).

\textbf{(2) Representation-dependent carriers.} If $Q_D$ depends on a privileged encoding, coordinate, or redescription choice, then admissibility criteria vary under admissible redescription,
contradicting Axiom~\ref{ax:miax}.

\textbf{(3) Representation-invariant structural carriers.} Suppose $Q_D$ is claimed to be purely structural and representation-invariant.
Then it is a candidate gate predicate. By Lemma~\ref{lem:unique-gate}, admissibility classification in a fixed domain is unique up to extensional equivalence of the gate.
Therefore either:
\begin{itemize}[leftmargin=2.2em]
\item $Q_D$ coincides extensionally with the gate predicate $G_D$ (in which case it is not an additional primitive; it is merely the gate under a new name), or
\item $Q_D$ disagrees with $G_D$ on some act, which would introduce a competing admissibility classification and is forbidden by Lemma~\ref{lem:unique-gate}.
\end{itemize}

In all cases, there is no primitive authorizer \emph{distinct} from the gate. Any admissible occurrence of $Q_D$ is therefore downstream bookkeeping inside an already gated domain, with no authority over eligibility.

\end{proof}

\begin{corollary}[Probability and typicality are not primitive]
\label{cor:prob}
No probabilistic typicality, measure dominance, Bayesian ranking, or thresholded confidence score can generate, repair, or authorize standing.
\end{corollary}

\begin{proof}
Any such scheme is a carrier attempt of type (1) in Theorem~\ref{thm:no-carriers}, hence inadmissible as primitive authorization.
\end{proof}

\begin{corollary}[Explanation cannot be the primitive]
\label{cor:explain}
Explanation, unification, compression, or coherence cannot function as primitive authorizers of standing.
\end{corollary}

\begin{proof}
If ``explanation'' is invoked to decide eligibility, it is either (i) a scalar/ordering surrogate (carrier attempt), or (ii) a scope/representation-dependent narrative (violating Axiom~\ref{ax:miax}),
or (iii) a post-hoc transformation of an already evaluated act (forbidden by Theorem~\ref{thm:no-repairs}). Therefore explanation cannot be primitive authorization.
\end{proof}

\begin{corollary}[Legitimacy conservation (closed regime)]
\label{cor:conservation}
Within the axioms, standing is conserved: no internal operation creates, amplifies, transfers, or recovers standing.
\end{corollary}

\begin{proof}
Creation/recovery would require a generator or repair, excluded by Theorems~\ref{thm:no-generators} and \ref{thm:no-repairs}.
Amplification or transfer by derived criteria would require a carrier attempt, excluded by Theorem~\ref{thm:no-carriers}.
\end{proof}

\section{Scope Discipline and Cross-Domain Transfer}
\label{sec:transfer}

\begin{lemma}[No illicit scope transport]
\label{lem:no-scope-transport}
Under Axiom~\ref{ax:scope}, no post-hoc change in domain typing may be used to legalize an otherwise-illicit act.
\end{lemma}

\begin{proof}
Axiom~\ref{ax:scope} states exactly this: an envelope shift used as repair is inadmissible. \end{proof}

\begin{definition}[Transfer claim]
A \emph{transfer claim} asserts that standing in a target domain $D_2$ is licensed on the basis of standing in a distinct source domain $D_1$,
without independent instantiation of target gating and target DAC.
\end{definition}

\begin{theorem}[Transfer without standing]
\label{thm:tws}
Let $D_1$ and $D_2$ be distinct domains. Standing does not inherit from $D_1$ to $D_2$ by default.
Standing in $D_2$ is defined only if $D_2$ independently instantiates a gate predicate and satisfies $\mathrm{DAC}_2$.
Any attempt to import standing from $D_1$ without independent target gating is illicit scope transport or carrier laundering and is fail-closed.
\end{theorem}

\begin{proof}
Assume a transfer claim licenses standing in $D_2$ from $D_1$ without independent target gating.

Either (i) it uses $D_1$'s gate (or admissibility/closure) to evaluate acts in $D_2$. This is a post-hoc envelope substitution, violating Lemma~\ref{lem:no-scope-transport}.
Or (ii) it uses a derived performance/similarity predicate to justify import. That predicate is a carrier attempt, forbidden as primitive authorization by Theorem~\ref{thm:no-carriers}.

Therefore, absent independent target gating and DAC, standing in $D_2$ is undefined (fail-closed).
\end{proof}

\begin{remark}[Regime translation is not standing lift]
Any ``translation'' or ``lift'' across regimes that is used to claim target standing without target gating is a transfer claim and is blocked by Theorem~\ref{thm:tws}.
\end{remark}

\section{Reference Fixation and Theory-Level Confirmation}
\label{sec:ref}

This section supplies the minimal standalone necessity claim linking standing to theory-level confirmation.

\begin{definition}[Admissible redescription]
An \emph{admissible redescription} is a transformation $d$ permitted by admissibility criteria $A$ within a domain.
\end{definition}

\begin{definition}[Fixation commutation]
Fixation \emph{commutes} with an admissible redescription $d$ if transporting a report under $d$ preserves what the report is about (no retargeting).
\end{definition}

\begin{definition}[Theory-level confirmation (minimal)]
A \emph{theory-level confirmation} claim is truth-apt only if the identity of the confirmed claim is invariant under admissible redescription.
\end{definition}

\begin{theorem}[Reference fixation is necessary for theory-level confirmation]
\label{thm:ref-fix}
If fixation fails to commute with admissible redescription, then theory-level confirmation is undefined. Predictive success alone cannot repair this failure.
\end{theorem}

\begin{proof}
If fixation does not commute, then admissible redescriptions change what the evidence is about. Then the ``same'' confirmation claim does not have a fixed target:
under admissible operations, different targets are being ``confirmed.'' This violates Axiom~\ref{ax:intell} (claim/evidence identity) or forces a representation-privileged choice,
contradicting Axiom~\ref{ax:miax}. Under Axiom~\ref{ax:biv}, there is no partial regime where such identity failure yields ``weakened'' confirmation; the confirmation relation is undefined.

Attempts to repair by success, probability, typicality, or any ranking surrogate are carrier attempts (Theorem~\ref{thm:no-carriers}) and are inadmissible at the boundary (Axiom~\ref{ax:ametric}).
Thus reference fixation commutation is necessary.
\end{proof}

\begin{remark}[Existence is fixed at the application layer]
\label{rem:interior-existence}
By Remark~\ref{rem:app-anchor}, the present paper works only where the standing-governed admissible side is already instantiated.
The uniqueness theorem below therefore does not attempt a lower derivation of admissible-side existence; it identifies the unique interior/closure of that already-instantiated admissible side.
\end{remark}

\section{Uniqueness of the Admissible Interior (No Discrete Multiplicity)}
\label{sec:unique}

\begin{definition}[Admissible interior]
\label{def:interior}
In the standing-instantiated regime fixed here, an \emph{admissible interior} $I$ is the closure of the already-instantiated standing-positive admissible side under admissibility-preserving operations.
Equivalently, it is the standing-governed admissible class closed under admissibility-preserving continuation inside the regime, maximal in the induced sense that no further same-regime admissible act may be added without leaving that standing-governed closure.
Two interiors $I_1$ and $I_2$ are distinct only if there exists a witness act $a$ such that $a\in I_1$ and $a\notin I_2$ (or vice versa).
\end{definition}

\begin{definition}[Admissible distinguishability]
Two purported interiors $I_1$ and $I_2$ are \emph{admissibly distinguishable} iff there exists an admissible act/predicate that witnesses their difference.
If no admissible witness exists, the interiors are admissibly indistinguishable.
\end{definition}

\begin{lemma}[Discrete multiplicity elimination (gate-form)]
\label{lem:discrete-elim}
Let $I_1$ and $I_2$ be purported admissible interiors (maximal admissible closures) for the same standing-governed regime.
Then $I_1=I_2$. In particular, discrete (categorical) plurality of admissible interiors is impossible.
\end{lemma}

\begin{proof}
Assume for contradiction that $I_1\neq I_2$. By extensionality of set difference there exists an act $a$ such that $a\in I_1$ and $a\notin I_2$ (or vice versa).
Define two candidate admissibility classifiers (gate predicates) for the same typed regime:
\[
G_1(x)\;:\!\!\iff\; x\in I_1,\qquad G_2(x)\;:\!\!\iff\; x\in I_2.
\]
By construction, $G_1(a)\neq G_2(a)$.

But Lemma~\ref{lem:unique-gate} states that within the intended regime (Axioms~\ref{ax:intell}, \ref{ax:ametric}, \ref{ax:biv}) there cannot exist two competing admissibility
classifications for the same domain/regime that disagree on any act: disagreement forces either indeterminate standing (violating intelligibility) or an admissibility-relevant selector
at the boundary (contradicting binary non-eventfulness and Lemma~\ref{lem:no-parameter}). Therefore $G_1$ and $G_2$ must coincide extensionally, contradicting $G_1(a)\neq G_2(a)$.
Hence $I_1=I_2$.
\end{proof}

\begin{theorem}[Uniqueness of the admissible interior]
\label{thm:unique}
In the standing-instantiated application regime fixed here, exactly one admissible interior exists.
\end{theorem}

\begin{proof}
By Remark~\ref{rem:app-anchor}, the present paper works only where the standing-governed admissible side is already instantiated.
In particular, the standing-positive admissible set is nonempty in the application regime under study, since non-degenerate reasoning is already operative there and standing is already instantiated at the upstream kernel layer.
Definition~\ref{def:interior} then names the closure of that already-instantiated standing-positive admissible side under admissibility-preserving operations, so an admissible interior exists in the regime under study.
If there were two distinct admissible interiors, Lemma~\ref{lem:discrete-elim} would be violated.
Therefore exactly one admissible interior exists.
\end{proof}

\section{Apex Standalone Closure Theorem}
\label{sec:apex}

\begin{theorem}[Apex Standalone Closure Theorem]
\label{thm:apex}
Under the standing-instantiated application regime of Remark~\ref{rem:app-anchor} and Axioms~\ref{ax:intell}--\ref{ax:biv}, the standing--admissibility regime is closed and unique:
\begin{enumerate}[label=(\roman*), leftmargin=2.2em]
\item \textbf{Closed standing regime:} standing is bivalent when defined (Lemma~\ref{lem:binary-standing}), gate-existence is forced (Lemma~\ref{lem:gate-existence}),
and there are no repairs (Theorem~\ref{thm:no-repairs}), no generators (Theorem~\ref{thm:no-generators}), and no primitive carriers (Theorem~\ref{thm:no-carriers}).
\item \textbf{Scope discipline:} no illicit scope transport (Lemma~\ref{lem:no-scope-transport}); no standing transfer without target gating and target DAC (Theorem~\ref{thm:tws}).
\item \textbf{Confirmation discipline:} theory-level confirmation requires reference fixation commutation; otherwise confirmation is undefined (Theorem~\ref{thm:ref-fix}).
\item \textbf{Uniqueness:} exactly one admissible interior exists, including elimination of discrete plural interiors (Theorem~\ref{thm:unique}).
\end{enumerate}
\end{theorem}

\begin{proof}
Items (i)--(iii) follow from the cited results. Item (iv) follows from Theorem~\ref{thm:unique}.
\end{proof}

\section{Degeneracy-Block Matrix (Proof Pointers)}
\label{sec:matrix}

\subsection{Matrix}
\label{subsec:matrix}

Each degeneracy class is blocked by at least one lemma/theorem proved above. The ``Blocked by'' column is therefore a proof pointer.

\setlength{\LTpre}{0pt}
\setlength{\LTpost}{0pt}
\begin{longtable}{>{\raggedright\arraybackslash}p{0.10\linewidth} >{\raggedright\arraybackslash}p{0.52\linewidth} >{\raggedright\arraybackslash}p{0.33\linewidth}}
\toprule
\textbf{ID} & \textbf{Degeneracy class (failure mode)} & \textbf{Blocked by (this paper)} \\
\midrule
\endfirsthead
\toprule
\textbf{ID} & \textbf{Degeneracy class (failure mode)} & \textbf{Blocked by (this paper)} \\
\midrule
\endhead

D1 & Graded standing / partial repair / ``almost legitimate.'' & Axiom~\ref{ax:biv}; Lemma~\ref{lem:binary-standing} \\
D2 & Metric/proximity transmissivity into gating; ``near-admissible.'' & Axiom~\ref{ax:ametric}; Lemma~\ref{lem:no-parameter}; Theorem~\ref{thm:no-carriers} \\
D3 & Probability/typicality as primitive authorization. & Corollary~\ref{cor:prob}; Theorem~\ref{thm:no-carriers} \\
D4 & Thresholded carrier authorization ($q>\tau$) or any scalar ranking surrogate. & Theorem~\ref{thm:no-carriers} \\
D5 & ``Transformational-invariant'' distribution preservation used as authorization. & Theorem~\ref{thm:no-carriers}; Axiom~\ref{ax:miax} \\
D6 & Post-hoc scope transport repair (narrow/expand/reframe to legalize failure). & Axiom~\ref{ax:scope}; Lemma~\ref{lem:no-scope-transport} \\
D7 & Cross-domain inheritance of standing (benchmark-to-world upgrades). & Theorem~\ref{thm:tws} \\
D8 & Standing accumulation/amplification (repetition/consensus/explanation/unification increases standing). & Theorem~\ref{thm:no-generators}; Theorem~\ref{thm:no-repairs}; Corollary~\ref{cor:explain} \\
D9 & Repair-by-explanation (explanation as primitive authorizer). & Corollary~\ref{cor:explain}; Theorem~\ref{thm:no-carriers} \\
D10 & Repair-by-unification/embedding (standing lift by consolidation). & Lemma~\ref{lem:no-scope-transport}; Theorem~\ref{thm:tws} \\
D11 & Silent redescription / privileged representation criterion. & Axiom~\ref{ax:miax}; Theorem~\ref{thm:no-carriers} \\
D12 & Cross-fixation relational transport or global structure from local articulation (scope smuggling). & Lemma~\ref{lem:no-scope-transport}; Definition of domain individuation \\
D13 & Boundary mutation / treating boundary descriptors as generative. & Axiom~\ref{ax:ametric}; Lemma~\ref{lem:no-parameter} \\
D14 & Continuous-parameter laundering (parameters become standing carriers). & Lemma~\ref{lem:no-parameter}; Theorem~\ref{thm:no-carriers} \\
D15 & Locality-violating ``extensions'' rescued by scope repair. & Lemma~\ref{lem:no-scope-transport} \\
D16 & Confirmation without reference fixation; predictive success treated as evidence absent fixed referent. & Theorem~\ref{thm:ref-fix} \\
D17 & Ordering reversal treated as weakened confirmation. & Axiom~\ref{ax:biv}; Theorem~\ref{thm:ref-fix} \\
D18 & Standing-collapse equivocation (confirmation failure vs standing loss). & Definition~\ref{def:standing}; Theorem~\ref{thm:ref-fix} \\
D19 & Identity drift misused as license for reference ambiguity. & Theorem~\ref{thm:ref-fix}; Axiom~\ref{ax:miax} \\
D20 & Non-transmissive recognition shortcut (authority by proposition alone). & Theorem~\ref{thm:tws}; Lemma~\ref{lem:no-scope-transport} \\
D21 & Regime translation misread as standing lift. & Remark in Section~\ref{sec:transfer}; Theorem~\ref{thm:tws} \\
D22 & Anchor vacuum / definitional admissibility triviality. & Lemma~\ref{lem:gate-existence}; Axiom~\ref{ax:intell} \\
D23 & Discrete plural admissible interiors (categorical multiplicity). & Lemma~\ref{lem:discrete-elim}; Theorem~\ref{thm:unique} \\
\bottomrule
\end{longtable}

\appendix
\section{Degeneracy Proof Appendix (By Cases)}
\label{app:deg}

This appendix provides a short proof for each degeneracy class D1--D23, showing it violates at least one axiom or proved theorem.

\begin{proposition}[D1: Graded standing]
Any graded-standing regime contradicts Axiom~\ref{ax:biv}.
\end{proposition}
\begin{proof}
Axiom~\ref{ax:biv} excludes ``almost'' legitimacy. Therefore any graded regime violates the axiom. \end{proof}

\begin{proposition}[D2: Metric transmissivity]
Any metric/proximity grounding of gating contradicts Axiom~\ref{ax:ametric} and Theorem~\ref{thm:no-carriers}.
\end{proposition}
\begin{proof}
Such grounding supplies an admissibility-relevant graded differentiator at the boundary, forbidden by Axiom~\ref{ax:ametric}; it is a carrier attempt blocked by Theorem~\ref{thm:no-carriers}. \end{proof}

\begin{proposition}[D3: Probability primitive]
Probability/typicality as primitive authorization is inadmissible.
\end{proposition}
\begin{proof}
Immediate from Corollary~\ref{cor:prob}. \end{proof}

\begin{proposition}[D4: Thresholded carriers]
Any $q>\tau$ authorization is a carrier attempt and is blocked.
\end{proposition}
\begin{proof}
It is a carrier attempt by definition and is excluded by Theorem~\ref{thm:no-carriers}. \end{proof}

\begin{proposition}[D5: Distribution-preservation authorization]
Any ``admissible iff preserves distribution'' rule is a carrier attempt and violates MIAX.
\end{proposition}
\begin{proof}
It substitutes a derived predicate for gating (Theorem~\ref{thm:no-carriers}). If the distribution depends on representation, MIAX is violated. \end{proof}

\begin{proposition}[D6: Post-hoc scope transport]
Post-hoc narrowing/expansion used as repair is inadmissible.
\end{proposition}
\begin{proof}
Immediate from Axiom~\ref{ax:scope} and Lemma~\ref{lem:no-scope-transport}. \end{proof}

\begin{proposition}[D7: Cross-domain inheritance]
Standing does not inherit across distinct domains without target gating and DAC.
\end{proposition}
\begin{proof}
Immediate from Theorem~\ref{thm:tws}. \end{proof}

\begin{proposition}[D8: Accumulation/amplification]
No repetition/consensus/explanation can amplify standing.
\end{proposition}
\begin{proof}
Amplification would create new legitimacy from old; this is excluded by Theorems~\ref{thm:no-generators}, \ref{thm:no-repairs} and Corollary~\ref{cor:explain}. \end{proof}

\begin{proposition}[D9: Explanation as repair]
Explanation cannot repair illegitimacy or authorize standing.
\end{proposition}
\begin{proof}
Immediate from Corollary~\ref{cor:explain} and Theorem~\ref{thm:no-repairs}. \end{proof}

\begin{proposition}[D10: Unification/embedding repair]
Unification used as standing lift is illicit.
\end{proposition}
\begin{proof}
It is either scope transport (Lemma~\ref{lem:no-scope-transport}) or a transfer claim without target gating (Theorem~\ref{thm:tws}). \end{proof}

\begin{proposition}[D11: Privileged representation]
Any privileged-encoding criterion violates MIAX.
\end{proposition}
\begin{proof}
Immediate from Axiom~\ref{ax:miax}. \end{proof}

\begin{proposition}[D12: Cross-fixation transport]
Cross-fixation comparability by silent envelope shift is scope transport.
\end{proposition}
\begin{proof}
If the comparison requires altering $F$ or $C$ without retyping the domain, it is illicit scope transport (Lemma~\ref{lem:no-scope-transport}). \end{proof}

\begin{proposition}[D13: Boundary mutation]
Treating boundary descriptors as generative violates binary AMetric boundary.
\end{proposition}
\begin{proof}
It makes boundary structure authoritative for gating, contradicting Axiom~\ref{ax:ametric} and Lemma~\ref{lem:no-parameter}. \end{proof}

\begin{proposition}[D14: Continuous-parameter laundering]
Continuous parameters cannot be used as standing carriers.
\end{proposition}
\begin{proof}
Any such use introduces admissibility-relevant parameterization (Lemma~\ref{lem:no-parameter}) and is a carrier attempt (Theorem~\ref{thm:no-carriers}). \end{proof}

\begin{proposition}[D15: Scope-rescued extensions]
Extensions rescued by scope repair are inadmissible.
\end{proposition}
\begin{proof}
They are post-hoc envelope shifts, excluded by Lemma~\ref{lem:no-scope-transport}. \end{proof}

\begin{proposition}[D16: Confirmation without fixation]
Theory-level confirmation is undefined without fixation commutation.
\end{proposition}
\begin{proof}
Immediate from Theorem~\ref{thm:ref-fix}. \end{proof}

\begin{proposition}[D17: ``Weakened'' confirmation]
Non-invariant confirmation is undefined, not weakened.
\end{proposition}
\begin{proof}
Axiom~\ref{ax:biv} excludes partial confirmation at this layer; Theorem~\ref{thm:ref-fix} yields undefinedness under fixation failure. \end{proof}

\begin{proposition}[D18: Standing/confirmation equivocation]
Standing undefinedness and illegitimacy are distinct and cannot be conflated.
\end{proposition}
\begin{proof}
By Definition~\ref{def:standing}, undefinedness is non-instantiation, not a graded value. Theorem~\ref{thm:ref-fix} shows confirmation may be undefined without asserting any standing value change. \end{proof}

\begin{proposition}[D19: Identity drift misuse]
Identity drift cannot license reference ambiguity for confirmation.
\end{proposition}
\begin{proof}
Using drift to permit retargeting violates MIAX and triggers the fixation non-commutation condition of Theorem~\ref{thm:ref-fix}. \end{proof}

\begin{proposition}[D20: Recognition shortcut]
Authority cannot be transmitted by proposition alone without target gating.
\end{proposition}
\begin{proof}
It is a transfer claim in disguise; blocked by Theorem~\ref{thm:tws}. \end{proof}

\begin{proposition}[D21: Regime translation as standing lift]
Translation is not authorization; it cannot lift standing.
\end{proposition}
\begin{proof}
If used to claim standing without target gating, it is a transfer claim blocked by Theorem~\ref{thm:tws}. \end{proof}

\begin{proposition}[D22: Anchor vacuum]
If no gate is instantiated, standing-bearing discourse collapses.
\end{proposition}
\begin{proof}
Lemma~\ref{lem:gate-existence} forces gate existence under intelligibility and irreversibility; otherwise non-degenerate discourse fails. \end{proof}

\begin{proposition}[D23: Discrete plural admissible interiors]
Discrete multiplicity of admissible interiors is impossible.
\end{proposition}
\begin{proof}
Immediate from Lemma~\ref{lem:discrete-elim} and Theorem~\ref{thm:unique}. \end{proof}

\section{Axiom Minimality and Independence Appendix}
\label{app:minindep}

This appendix has two functions.

\begin{enumerate}[label=(\alph*), leftmargin=2.2em]
\item \textbf{Minimality.} Identify which axioms are truly required for each lemma/theorem in the main text (including transitive dependencies through cited results).
\item \textbf{Independence.} Provide witness constructions showing that, for each axiom, dropping it while keeping the others admits at least one degeneracy class (D1--D23) or breaks at least one derived theorem.
\end{enumerate}

\subsection{Axiom--Result Dependency Matrix}
\label{app:dep}

The application-regime anchor of Remark~\ref{rem:app-anchor} is an upstream scope condition rather than an additional internal axiom of the present paper.
Accordingly, the matrix below tracks only the internal axioms of Section~\ref{sec:axioms}; whenever existence of the admissible side is needed, that requirement is treated as fixed at application rather than encoded as an extra matrix column.

We use the shorthand:
\[
\mathbf{I}=\text{Intelligibility (Axiom~\ref{ax:intell})},\;
\mathbf{R}=\text{Irreversibility (Axiom~\ref{ax:irr})},\;
\mathbf{C}=\text{Composition (Axiom~\ref{ax:comp})},\;
\mathbf{S}=\text{Scope discipline (Axiom~\ref{ax:scope})},\;
\mathbf{M}=\text{MIAX (Axiom~\ref{ax:miax})},\;
\mathbf{A}=\text{AMetric boundary (Axiom~\ref{ax:ametric})},\;
\mathbf{B}=\text{Bivalence (Axiom~\ref{ax:biv})}.
\]
Kernel instantiation / non-vacuity is inherited from the application-regime note in Section~\ref{sec:axioms} rather than introduced here as a local axiom.
The matrix below therefore tracks only \emph{local axiom use} inside this standalone consequence derivation.
A checkmark indicates that the axiom is required by the result \emph{as a consequence of its proof chain} (i.e.\ including dependencies on cited lemmas/theorems).

\newcommand{\X}{\checkmark}

\setlength{\LTpre}{0pt}
\setlength{\LTpost}{0pt}
\begin{longtable}{>{\raggedright\arraybackslash}p{0.33\linewidth}ccccccc}
\toprule
\textbf{Result} & $\mathbf{I}$ & $\mathbf{R}$ & $\mathbf{C}$ & $\mathbf{S}$ & $\mathbf{M}$ & $\mathbf{A}$ & $\mathbf{B}$\\
\midrule
\endfirsthead
\toprule
\textbf{Result} & $\mathbf{I}$ & $\mathbf{R}$ & $\mathbf{C}$ & $\mathbf{S}$ & $\mathbf{M}$ & $\mathbf{A}$ & $\mathbf{B}$\\
\midrule
\endhead

Lemma~\ref{lem:gate-existence} (gate existence) & \X & \X &  &  &  &  &  \\
Lemma~\ref{lem:binary-standing} (binary standing) &  &  &  &  &  &  & \X \\
Lemma~\ref{lem:no-parameter} (no parameterization) &  &  &  &  &  & \X &  \\
Lemma~\ref{lem:unique-gate} (no competing gates) & \X &  &  &  &  & \X & \X \\
Theorem~\ref{thm:no-repairs} (no repairs) &  & \X &  &  &  &  &  \\
Theorem~\ref{thm:no-generators} (no generators) &  & \X &  &  &  &  &  \\
Lemma~\ref{lem:authorizer-is-gate} (authorizer $\Rightarrow$ gate) &  &  &  &  &  &  & \X \\
Theorem~\ref{thm:no-carriers} (no carriers) & \X &  &  &  & \X & \X & \X \\
Corollary~\ref{cor:prob} (probability not primitive) & \X &  &  &  & \X & \X & \X \\
Corollary~\ref{cor:explain} (explanation not primitive) & \X & \X &  &  & \X & \X & \X \\
Corollary~\ref{cor:conservation} (conservation) & \X & \X &  &  & \X & \X & \X \\
Lemma~\ref{lem:no-scope-transport} (no scope transport) &  &  &  & \X &  &  &  \\
Theorem~\ref{thm:tws} (transfer without standing) & \X &  &  & \X & \X & \X & \X \\
Theorem~\ref{thm:ref-fix} (fixation necessary) & \X &  &  &  & \X & \X & \X \\
Lemma~\ref{lem:discrete-elim} (no discrete interiors) & \X &  &  &  &  & \X & \X \\
Theorem~\ref{thm:unique} (uniqueness) & \X &  &  &  &  & \X & \X \\
Theorem~\ref{thm:apex} (apex closure) & \X & \X &  & \X & \X & \X & \X \\
\bottomrule
\end{longtable}

\begin{remark}[About the composition axiom]
As packaged here, Axiom~\ref{ax:comp} is not required in any proof chain because MIAX (Axiom~\ref{ax:miax}) is assumed as primitive.
In the corpus, composition plays a deeper role: it is part of the route by which MIAX is forced rather than assumed.
Accordingly, \emph{for this standalone derivation}, $\mathbf{C}$ is non-minimal; see Proposition~\ref{prop:comp-nonminimal}.
\end{remark}

\begin{proposition}[Composition is non-minimal in the standalone derivation]
\label{prop:comp-nonminimal}
In the standalone proof architecture of this paper, Axiom~\ref{ax:comp} is not required by any derived result.
\end{proposition}

\begin{proof}
Inspection of the dependency matrix shows that no lemma, theorem, or corollary proved in the main text depends on Axiom~\ref{ax:comp}. In this standalone package, MIAX is assumed directly as Axiom~\ref{ax:miax}; composition therefore does not enter any proof chain, even though it remains part of the broader corpus route by which MIAX is forced.
\end{proof}

\subsection{Minimal axiom sets per result}
\label{app:minsets}

For each derived result, we record a minimal \emph{internal} axiom subset sufficient for the proof chain once the standing-instantiated application regime of Remark~\ref{rem:app-anchor} is fixed.

\begin{itemize}[leftmargin=2.2em]
\item Lemma~\ref{lem:gate-existence}: $\{\mathbf{I},\mathbf{R}\}$.
\item Lemma~\ref{lem:binary-standing}: $\{\mathbf{B}\}$.
\item Lemma~\ref{lem:no-parameter}: $\{\mathbf{A}\}$.
\item Lemma~\ref{lem:unique-gate}: $\{\mathbf{I},\mathbf{A},\mathbf{B}\}$.
\item Theorem~\ref{thm:no-repairs}: $\{\mathbf{R}\}$ (with standing defined).
\item Theorem~\ref{thm:no-generators}: $\{\mathbf{R}\}$ via Theorem~\ref{thm:no-repairs}.
\item Lemma~\ref{lem:authorizer-is-gate}: $\{\mathbf{B}\}$.
\item Theorem~\ref{thm:no-carriers}: $\{\mathbf{I},\mathbf{M},\mathbf{A},\mathbf{B}\}$ plus Lemmas~\ref{lem:no-parameter} and \ref{lem:unique-gate}.
\item Corollary~\ref{cor:prob}: Theorem~\ref{thm:no-carriers} (hence $\{\mathbf{I},\mathbf{M},\mathbf{A},\mathbf{B}\}$).
\item Corollary~\ref{cor:explain}: Theorem~\ref{thm:no-carriers} plus Theorem~\ref{thm:no-repairs} (hence $\{\mathbf{I},\mathbf{R},\mathbf{M},\mathbf{A},\mathbf{B}\}$).
\item Corollary~\ref{cor:conservation}: Theorems~\ref{thm:no-repairs}, \ref{thm:no-generators}, \ref{thm:no-carriers} (hence $\{\mathbf{I},\mathbf{R},\mathbf{M},\mathbf{A},\mathbf{B}\}$).
\item Lemma~\ref{lem:no-scope-transport}: $\{\mathbf{S}\}$.
\item Theorem~\ref{thm:tws}: $\{\mathbf{S}\}$ plus Theorem~\ref{thm:no-carriers} (hence $\{\mathbf{I},\mathbf{S},\mathbf{M},\mathbf{A},\mathbf{B}\}$).
\item Theorem~\ref{thm:ref-fix}: $\{\mathbf{I},\mathbf{M},\mathbf{A},\mathbf{B}\}$ plus Theorem~\ref{thm:no-carriers}.
\item Lemma~\ref{lem:discrete-elim}: $\{\mathbf{I},\mathbf{A},\mathbf{B}\}$ via Lemma~\ref{lem:unique-gate} (no competing admissibility classifications).
\item Theorem~\ref{thm:unique}: application-regime anchor (Remark~\ref{rem:app-anchor}) plus Lemma~\ref{lem:discrete-elim}; internally, $\{\mathbf{I},\mathbf{A},\mathbf{B}\}$ suffice for the uniqueness part once existence is fixed at application.
\item Theorem~\ref{thm:apex}: application-regime anchor plus $\{\mathbf{I},\mathbf{R},\mathbf{S},\mathbf{M},\mathbf{A},\mathbf{B}\}$.
\end{itemize}

\subsection{Independence models in a fixed signature (explicit \texorpdfstring{$\Sigma$}{Sigma}-structures)}
\label{app:indep}

\paragraph{What is proved here.}
We define a \emph{single} finite parametric family of first-order $\Sigma$-structures $\mathcal{M}(\theta)$ indexed by a one-place ``toggle'' parameter
\[
\theta\in \Theta:=\{\varnothing,\mathbf{I},\mathbf{R},\mathbf{C},\mathbf{S},\mathbf{M},\mathbf{A},\mathbf{B}\}.
\]
The base setting $\theta=\varnothing$ satisfies $T_0$ and all formal axiom cores $(X_\Sigma)$.
For each nonempty $X\in\Theta$, the specialization $\mathcal{M}(X)$ satisfies $T_0$ and all cores except $(X_\Sigma)$, and in addition satisfies a designated
formal degeneracy formula $(D^X_\Sigma)$ witnessing a reopened failure mode.
Thus independence is shown \emph{in the intended nontrivial class} $T_0$ by \emph{controlled perturbations} of one fixed base model (Section~\ref{app:paramfam}),
not by vacuous empty-sort models and not by unrelated hand-built structures.

\paragraph{Adequacy note.}
The schemata $(X_\Sigma)$ below are \emph{formal cores} of the corresponding informal axioms in the main text: they isolate precisely the structure actually used in the closure proofs.
Appendix~\ref{app:indep} therefore establishes independence of these cores. (No claim is made that the full informal axioms are \emph{equivalent} to their cores.)

\subsubsection{Common signature \texorpdfstring{$\Sigma$}{Sigma} with functions}

We use one fixed first-order signature $\Sigma$ with sort-predicates (unary predicates) and with \emph{functions} for standing value and metric score:
\begin{itemize}[leftmargin=2.2em]
\item Sort-predicates: $\mathrm{Dom}(x)$, $\mathrm{Act}(x)$, $\mathrm{Tar}(x)$, $\mathrm{Val}(x)$, $\mathrm{Ev}(x)$.
\item Constants:
\[
d_1,d_2;\quad a_0,a_1,a_2,a_3;\quad e_0,e_1,e_2;\quad t_0,t_1;\quad v_0,v_1,v_h,\tau.
\]
\item Relations:
  \begin{itemize}[leftmargin=2.2em]
  \item $\mathrm{Eval}(d,a)$: act $a$ is evaluated in domain $d$.
  \item $\mathrm{Red}(x,y)$: $y$ is an admissible redescription of $x$ (used for both evidence and acts).
  \item $\mathrm{Same}(a,b)$: $a$ and $b$ are treated as the ``same act'' (identity-preserving continuation).
  \item $\mathrm{Sub}(d',d)$: $d'$ is a post-hoc narrowing (sub-envelope) of $d$.
  \item $\mathrm{About}(e,t)$: evidence-act $e$ is about target $t$.
  \item $\mathrm{LT}(v,w)$: ordering on values (used to express threshold rules).
  \end{itemize}
\item Functions:
  \begin{itemize}[leftmargin=2.2em]
  \item $sv(d,a)$: the standing value assigned to act $a$ in domain $d$.
  \item $score(a)$: a metric/probabilistic score assigned to act $a$.
  \end{itemize}
\end{itemize}

\subsubsection{Background nontriviality theory \texorpdfstring{$T_0$}{T0}}
\label{app:T0}

$T_0$ forces the signature to be non-degenerate and ensures the axioms are not satisfied vacuously:
\begin{itemize}[leftmargin=2.2em]
\item \textbf{Typed constants and distinctness:}
\[
\mathrm{Dom}(d_1)\wedge \mathrm{Dom}(d_2)\wedge d_1\neq d_2,
\]
\[
\bigwedge_{i=0}^3 \mathrm{Act}(a_i)\;\wedge\; \bigwedge_{0\le i<j\le 3} a_i\neq a_j,
\]
\[
\bigwedge_{i=0}^2 \mathrm{Ev}(e_i)\;\wedge\; \bigwedge_{0\le i<j\le 2} e_i\neq e_j,
\]
\[
\mathrm{Tar}(t_0)\wedge \mathrm{Tar}(t_1)\wedge t_0\neq t_1,
\]
\[
\mathrm{Val}(v_0)\wedge \mathrm{Val}(v_1)\wedge \mathrm{Val}(v_h)\wedge \mathrm{Val}(\tau)\wedge (v_0\neq v_1)\wedge (v_h\neq v_0)\wedge (v_h\neq v_1).
\]
\item \textbf{Nontrivial evaluation and scope:} $\mathrm{Sub}(d_2,d_1)$ and $\mathrm{Eval}(d_k,a_i)$ for all $k\in\{1,2\}$ and all $i\in\{0,1,2,3\}$.
\item \textbf{Nontrivial redescription chains:} $\mathrm{Red}(e_0,e_1)\wedge \mathrm{Red}(e_1,e_2)$ and $\mathrm{Red}(a_0,a_1)\wedge \mathrm{Red}(a_1,a_2)$.
\item \textbf{Nontrivial order:} $\mathrm{LT}(v_0,\tau)\wedge \mathrm{LT}(\tau,v_1)$.
\end{itemize}

\subsubsection{Formal axiom cores in \texorpdfstring{$\Sigma$}{Sigma}}

\begin{itemize}[leftmargin=2.2em]
\item \textbf{(I$_\Sigma$) Intelligibility core.} Fixation commutes under admissible redescription for evidence:
\[
\forall e,e',t\;\big(\mathrm{Ev}(e)\wedge \mathrm{Ev}(e')\wedge \mathrm{Red}(e,e') \rightarrow (\mathrm{About}(e,t)\leftrightarrow \mathrm{About}(e',t))\big).
\]
\item \textbf{(R$_\Sigma$) Irreversibility core.} Standing is invariant under ``same act'' continuation:
\[
\forall d,a,a'\;\big(\mathrm{Dom}(d)\wedge \mathrm{Act}(a)\wedge \mathrm{Act}(a')\wedge \mathrm{Eval}(d,a)\wedge \mathrm{Eval}(d,a')\wedge \mathrm{Same}(a,a')\rightarrow sv(d,a)=sv(d,a')\big).
\]
\item \textbf{(C$_\Sigma$) Composition core.} Admissible redescription composes (transitivity):
\[
\forall x,y,z\;(\mathrm{Red}(x,y)\wedge \mathrm{Red}(y,z)\rightarrow \mathrm{Red}(x,z)).
\]
\item \textbf{(S$_\Sigma$) Scope discipline core.} No post-hoc narrowing can legalize (or otherwise change) standing:
\[
\forall d,d',a\;\big(\mathrm{Dom}(d)\wedge \mathrm{Dom}(d')\wedge \mathrm{Act}(a)\wedge \mathrm{Sub}(d',d)\wedge \mathrm{Eval}(d,a)\wedge \mathrm{Eval}(d',a)\rightarrow sv(d',a)=sv(d,a)\big).
\]
\item \textbf{(M$_\Sigma$) MIAX core.} Standing is invariant under admissible redescription of acts:
\[
\forall d,a,a'\;\big(\mathrm{Dom}(d)\wedge \mathrm{Act}(a)\wedge \mathrm{Act}(a')\wedge \mathrm{Eval}(d,a)\wedge \mathrm{Eval}(d,a')\wedge \mathrm{Red}(a,a')\rightarrow sv(d,a)=sv(d,a')\big).
\]
\item \textbf{(A$_\Sigma$) AMetric core.} The score structure is non-discriminative (constant on acts):
\[
\forall a,b\;\big(\mathrm{Act}(a)\wedge \mathrm{Act}(b)\rightarrow score(a)=score(b)\big).
\]
\item \textbf{(B$_\Sigma$) Bivalence core.} Only $v_0$ and $v_1$ occur as standing values for evaluated acts:
\[
\forall d,a\;\big(\mathrm{Dom}(d)\wedge \mathrm{Act}(a)\wedge \mathrm{Eval}(d,a)\rightarrow (sv(d,a)=v_0\vee sv(d,a)=v_1)\big).
\]
\end{itemize}

\subsubsection{Formal degeneracy formulas used in the independence claims}

We use the following $\Sigma$-formulas as ``reopened degeneracies'' (one per axiom-drop witness):

\begin{itemize}[leftmargin=2.2em]
\item \textbf{(D$^{\mathbf{I}}_\Sigma$)} Fixation drift (D16-type):
\[
\exists e,e'\;\big(\mathrm{Ev}(e)\wedge \mathrm{Ev}(e')\wedge \mathrm{Red}(e,e')\wedge \mathrm{About}(e,t_0)\wedge \mathrm{About}(e',t_1)\big).
\]
\item \textbf{(D$^{\mathbf{R}}_\Sigma$)} Same-act repair (D1/No-Repairs failure):
\[
\exists d\;\big(\mathrm{Dom}(d)\wedge \mathrm{Same}(a_2,a_3)\wedge \mathrm{Eval}(d,a_2)\wedge \mathrm{Eval}(d,a_3)\wedge sv(d,a_2)=v_0\wedge sv(d,a_3)=v_1\big).
\]
\item \textbf{(D$^{\mathbf{C}}_\Sigma$)} Non-composition of redescription:
\[
\mathrm{Red}(e_0,e_1)\wedge \mathrm{Red}(e_1,e_2)\wedge \neg \mathrm{Red}(e_0,e_2).
\]
\item \textbf{(D$^{\mathbf{S}}_\Sigma$)} Scope legalization (D6-type):
\[
\mathrm{Sub}(d_2,d_1)\wedge \mathrm{Eval}(d_1,a_3)\wedge \mathrm{Eval}(d_2,a_3)\wedge sv(d_1,a_3)=v_0\wedge sv(d_2,a_3)=v_1.
\]
\item \textbf{(D$^{\mathbf{M}}_\Sigma$)} Redescription-variant standing (D11-type):
\[
\exists d\;\big(\mathrm{Dom}(d)\wedge \mathrm{Red}(a_0,a_1)\wedge \mathrm{Eval}(d,a_0)\wedge \mathrm{Eval}(d,a_1)\wedge sv(d,a_0)=v_0\wedge sv(d,a_1)=v_1\big).
\]
\item \textbf{(D$^{\mathbf{A}}_\Sigma$)} Threshold carrier authorizes standing (D4-type):
\[
\exists d\;\Big(\mathrm{Dom}(d)\wedge \neg \mathrm{LT}(\tau,score(a_0))\wedge \mathrm{LT}(\tau,score(a_3))\wedge
\forall x\big(\mathrm{Act}(x)\wedge \mathrm{Eval}(d,x)\rightarrow (sv(d,x)=v_1\leftrightarrow \mathrm{LT}(\tau,score(x)))\big)\Big).
\]
\item \textbf{(D$^{\mathbf{B}}_\Sigma$)} Third standing value (graded/ternary standing):
\[
\exists d\;\big(\mathrm{Dom}(d)\wedge \mathrm{Eval}(d,a_3)\wedge sv(d,a_3)=v_h\big).
\]
\end{itemize}

\subsubsection{A base model \texorpdfstring{$\mathcal{M}_0$}{M0} satisfying \texorpdfstring{$T_0$}{T0} and all cores}

We first define one concrete finite structure $\mathcal{M}_0$ that satisfies $T_0$ and all $(X_\Sigma)$.
Let the universe be
\[
U=\{d_1,d_2,a_0,a_1,a_2,a_3,e_0,e_1,e_2,t_0,t_1,v_0,v_1,v_h,\tau\}.
\]
Interpret the sort predicates by:
\[
\mathrm{Dom}^{\mathcal{M}_0}=\{d_1,d_2\},\;
\mathrm{Act}^{\mathcal{M}_0}=\{a_0,a_1,a_2,a_3\},\;
\mathrm{Ev}^{\mathcal{M}_0}=\{e_0,e_1,e_2\},\;
\mathrm{Tar}^{\mathcal{M}_0}=\{t_0,t_1\},\;
\mathrm{Val}^{\mathcal{M}_0}=\{v_0,v_1,v_h,\tau\}.
\]
Set:
\[
\mathrm{Sub}^{\mathcal{M}_0}=\{(d_2,d_1)\},\qquad
\mathrm{Eval}^{\mathcal{M}_0}=\{(d_k,a_i)\;:\;k\in\{1,2\},\, i\in\{0,1,2,3\}\}.
\]
Let the redescription relation be the transitive closure of the required chains:
\[
\mathrm{Red}^{\mathcal{M}_0}=\{(e_0,e_1),(e_1,e_2),(e_0,e_2)\}\cup\{(a_0,a_1),(a_1,a_2),(a_0,a_2)\}.
\]
Let $\mathrm{Same}^{\mathcal{M}_0}=\emptyset$ (no nontrivial same-act identifications in the base model).
Let aboutness be fixation-preserving:
\[
\mathrm{About}^{\mathcal{M}_0}=\{(e_0,t_0),(e_1,t_0),(e_2,t_0)\}.
\]
Let the order relation include the nontrivial threshold facts:
\[
\mathrm{LT}^{\mathcal{M}_0}=\{(v_0,\tau),(\tau,v_1)\}.
\]
Finally define the functions:
\[
score^{\mathcal{M}_0}(a)=v_0\ \text{for all}\ a\in \mathrm{Act}^{\mathcal{M}_0},\qquad
sv^{\mathcal{M}_0}(d,a)=v_0\ \text{for all}\ d\in \mathrm{Dom}^{\mathcal{M}_0},\, a\in \mathrm{Act}^{\mathcal{M}_0}.
\]
It is immediate that $\mathcal{M}_0\models T_0$ and $\mathcal{M}_0\models \bigwedge_X (X_\Sigma)$.

\subsubsection{Uniform parametric witness family \texorpdfstring{$\mathcal{M}(\theta)$}{M(theta)} (controlled perturbations)}
\label{app:paramfam}

Rather than presenting seven unrelated hand-tuned countermodels, we package the entire independence appendix as a \emph{single} finite parametric
$\Sigma$-structure family $\mathcal{M}(\theta)$ with one ``toggle'' parameter $\theta$.
Each nontrivial setting of $\theta$ flips \emph{exactly one} axiom-core $(X_\Sigma)$ while preserving $T_0$ and all remaining cores.

\paragraph{Toggle set.}
Let
\[
\Theta:=\{\varnothing,\mathbf{I},\mathbf{R},\mathbf{C},\mathbf{S},\mathbf{M},\mathbf{A},\mathbf{B}\},
\]
where $\varnothing$ means ``no perturbation'' (the base model) and each nonempty element indicates \emph{which} core axiom is to be violated.

\paragraph{Definition of \texorpdfstring{$\mathcal{M}(\theta)$}{M(theta)}.}
Fix the universe, sorts, and the interpretations of $\mathrm{Sub}$, $\mathrm{Eval}$, and $\mathrm{LT}$ to be exactly those of $\mathcal{M}_0$.
All symbols not listed below are interpreted as in $\mathcal{M}_0$.
The varying symbols are $\mathrm{Red}$, $\mathrm{Same}$, $\mathrm{About}$, and the functions $score(\cdot)$ and $sv(\cdot,\cdot)$, defined by cases:

\begin{itemize}[leftmargin=2.2em]
\item \textbf{Redescription:}
\[
\mathrm{Red}^{\mathcal{M}(\theta)} :=
\begin{cases}
\mathrm{Red}^{\mathcal{M}_0}\setminus\{(e_0,e_2)\}, & \text{if }\theta=\mathbf{C},\\
\mathrm{Red}^{\mathcal{M}_0}, & \text{otherwise.}
\end{cases}
\]
\item \textbf{Same-act relation:}
\[
\mathrm{Same}^{\mathcal{M}(\theta)} :=
\begin{cases}
\{(a_2,a_3)\}, & \text{if }\theta=\mathbf{R},\\
\emptyset, & \text{otherwise.}
\end{cases}
\]
\item \textbf{Aboutness (fixation):}
\[
\mathrm{About}^{\mathcal{M}(\theta)} :=
\begin{cases}
\{(e_0,t_0),(e_1,t_0),(e_2,t_0),(e_2,t_1)\}, & \text{if }\theta=\mathbf{I},\\
\{(e_0,t_0),(e_1,t_0),(e_2,t_0)\}, & \text{otherwise.}
\end{cases}
\]
\item \textbf{Score function:}
\[
score^{\mathcal{M}(\theta)}(x):=
\begin{cases}
v_1, & \text{if }\theta=\mathbf{A}\ \text{and}\ x=a_3,\\
v_0, & \text{otherwise (in particular for all acts if }\theta\neq \mathbf{A}\text{).}
\end{cases}
\]
\item \textbf{Standing-value function:} for $d\in\mathrm{Dom}$ and $x\in\mathrm{Act}$,
\[
sv^{\mathcal{M}(\theta)}(d,x):=
\begin{cases}
v_1, & \text{if }\theta=\mathbf{A}\ \text{and}\ \mathrm{LT}(\tau,score^{\mathcal{M}(\theta)}(x)),\\
v_0, & \text{if }\theta=\mathbf{A}\ \text{and}\ \neg \mathrm{LT}(\tau,score^{\mathcal{M}(\theta)}(x)),\\[4pt]
v_h, & \text{if }\theta=\mathbf{B}\ \text{and}\ x=a_3,\\
v_1, & \text{if }\theta=\mathbf{R}\ \text{and}\ x=a_3,\\
v_1, & \text{if }\theta=\mathbf{S}\ \text{and}\ d=d_2\ \text{and}\ x=a_3,\\
v_1, & \text{if }\theta=\mathbf{M}\ \text{and}\ x=a_1,\\
v_0, & \text{otherwise.}
\end{cases}
\]
\end{itemize}

\paragraph{Uniform witness theorem.}
The family $\mathcal{M}(\theta)$ supplies, by controlled perturbation, all axiom-drop witnesses.

\begin{theorem}[Uniform independence by controlled perturbations]
\label{thm:uniform-indep}
For every $\theta\in\Theta$, $\mathcal{M}(\theta)\models T_0$. Moreover:
\begin{enumerate}[label=(\alph*), leftmargin=2.2em]
\item $\mathcal{M}(\varnothing)=\mathcal{M}_0$ and $\mathcal{M}(\varnothing)\models \bigwedge_{X\in\{\mathbf{I},\mathbf{R},\mathbf{C},\mathbf{S},\mathbf{M},\mathbf{A},\mathbf{B}\}} (X_\Sigma)$.
\item For each $X\in\{\mathbf{I},\mathbf{R},\mathbf{C},\mathbf{S},\mathbf{M},\mathbf{A},\mathbf{B}\}$:
\[
\mathcal{M}(X)\models T_0\ \wedge\ \bigwedge_{Y\neq X}(Y_\Sigma)\ \wedge\ \neg (X_\Sigma)\ \wedge\ (D^X_\Sigma).
\]
\end{enumerate}
\end{theorem}

\begin{proof}
$T_0$ is preserved for all $\theta$ because the interpretations of all symbols appearing in $T_0$ are either fixed as in $\mathcal{M}_0$
(sorts, $\mathrm{Dom}$/$\mathrm{Act}$/$\mathrm{Ev}$/$\mathrm{Tar}$/$\mathrm{Val}$, constants, $\mathrm{Sub}$, $\mathrm{Eval}$, the required chain facts
$\mathrm{Red}(e_0,e_1)$, $\mathrm{Red}(e_1,e_2)$, $\mathrm{Red}(a_0,a_1)$, $\mathrm{Red}(a_1,a_2)$, and $\mathrm{LT}(v_0,\tau)$, $\mathrm{LT}(\tau,v_1)$),
or are modified in a way that does not remove the required $T_0$ instances (the only modification to $\mathrm{Red}$ is removal of $(e_0,e_2)$, which is not required by $T_0$).

For part (a), this is by construction: $\varnothing$ triggers none of the case-clauses, so all interpretations coincide with $\mathcal{M}_0$.

For part (b), check each $X$:

\textbf{$X=\mathbf{I}$.} $\theta=\mathbf{I}$ changes only $\mathrm{About}$ so that $\mathrm{Red}(e_1,e_2)$ holds while $\mathrm{About}(e_1,t_0)$ and $\mathrm{About}(e_2,t_1)$, hence $(I_\Sigma)$ fails and $(D^{\mathbf{I}}_\Sigma)$ holds. All other cores remain true because $\mathrm{Same}$ is empty, $\mathrm{Red}$ is transitive, $sv$ is constant $v_0$, and $score$ is constant $v_0$.

\textbf{$X=\mathbf{R}$.} $\theta=\mathbf{R}$ changes only $\mathrm{Same}$ and $sv$ so that $\mathrm{Same}(a_2,a_3)$ and $sv(d,a_2)=v_0\neq v_1=sv(d,a_3)$, hence $(R_\Sigma)$ fails and $(D^{\mathbf{R}}_\Sigma)$ holds. MIAX holds because all redescription-related acts $a_0,a_1,a_2$ still have standing $v_0$; scope discipline holds because $sv$ is domain-invariant; the remaining cores are unchanged.

\textbf{$X=\mathbf{C}$.} $\theta=\mathbf{C}$ changes only $\mathrm{Red}$ by removing $(e_0,e_2)$ while retaining $\mathrm{Red}(e_0,e_1)$ and $\mathrm{Red}(e_1,e_2)$, hence $(C_\Sigma)$ fails and $(D^{\mathbf{C}}_\Sigma)$ holds. The other cores are unaffected since they quantify only over existing $\mathrm{Red}$ pairs (intelligibility remains true for the remaining $\mathrm{Red}$ links), and $sv$/$score$ are as in $\mathcal{M}_0$.

\textbf{$X=\mathbf{S}$.} $\theta=\mathbf{S}$ changes only $sv(d_2,a_3)$ to $v_1$ while keeping $sv(d_1,a_3)=v_0$, hence $(S_\Sigma)$ fails using $\mathrm{Sub}(d_2,d_1)$ and $(D^{\mathbf{S}}_\Sigma)$ holds. All other cores remain true: MIAX is preserved on the redescription chain $a_0\to a_1\to a_2$ (all remain $v_0$), bivalence holds, and $\mathrm{Same}=\emptyset$ keeps $(R_\Sigma)$ true.

\textbf{$X=\mathbf{M}$.} $\theta=\mathbf{M}$ changes only $sv(d,a_1)$ to $v_1$, while $\mathrm{Red}(a_0,a_1)$ holds and $sv(d,a_0)=v_0$, so $(M_\Sigma)$ fails and $(D^{\mathbf{M}}_\Sigma)$ holds. Scope discipline remains true because the standing change is domain-invariant; other cores are unchanged.

\textbf{$X=\mathbf{A}$.} $\theta=\mathbf{A}$ changes $score(a_3)$ to $v_1$ (others $v_0$) and sets standing by the threshold rule, so $(A_\Sigma)$ fails and $(D^{\mathbf{A}}_\Sigma)$ holds by definition. MIAX remains true because redescription-linked acts $a_0,a_1,a_2$ all have $score=v_0$ and hence $sv=v_0$; scope discipline holds because the threshold rule is domain-invariant; bivalence holds since $sv\in\{v_0,v_1\}$.

\textbf{$X=\mathbf{B}$.} $\theta=\mathbf{B}$ sets $sv(d,a_3)=v_h$, so $(B_\Sigma)$ fails and $(D^{\mathbf{B}}_\Sigma)$ holds. All other cores remain true since the only altered standing values are identical across domains (preserving scope discipline), do not interact with $\mathrm{Same}$ (empty), and do not violate MIAX on the redescription chain (which does not involve $a_3$ in $T_0$).
\end{proof}


\subsubsection{Minimal-change property in a Hamming metric}
\label{app:minchange}

\begin{definition}[Hamming distance on $\Sigma$-interpretations]
Fix the universe $U$ and constant/sort interpretations as in $\mathcal{M}_0$.
For two $\Sigma$-structures $\mathcal{N},\mathcal{N}'$ on $U$, define
\[
d(\mathcal{N},\mathcal{N}'):=\sum_{R\in\mathrm{Rel}(\Sigma)} \big|R^{\mathcal{N}}\triangle R^{\mathcal{N}'}\big|
\;+\;
\sum_{f\in\mathrm{Fun}(\Sigma)} \big|\{\bar u\in U^{\mathrm{arity}(f)}:\ f^{\mathcal{N}}(\bar u)\neq f^{\mathcal{N}'}(\bar u)\}\big|.
\]
This is the Hamming distance on relation-graphs and function-graphs (with symmetric difference $\triangle$ for relations).
\end{definition}

\begin{lemma}[Minimal-change toggles]
\label{lem:minchange}
Let $X\in\{\mathbf{I},\mathbf{R},\mathbf{C},\mathbf{S},\mathbf{M},\mathbf{A},\mathbf{B}\}$ and define:
\[
\begin{array}{c|ccccccc}
X & \mathbf{I} & \mathbf{R} & \mathbf{C} & \mathbf{S} & \mathbf{M} & \mathbf{A} & \mathbf{B}\\
\hline
k_X & 1 & 3 & 1 & 1 & 2 & 3 & 2
\end{array}
\]
Then:
\begin{enumerate}[label=(\alph*), leftmargin=2.2em]
\item $d(\mathcal{M}_0,\mathcal{M}(X))=k_X$.
\item For any $\Sigma$-structure $\mathcal{N}$ on $U$ satisfying
\[
\mathcal{N}\models T_0\ \wedge\ \bigwedge_{Y\neq X}(Y_\Sigma)\ \wedge\ \neg(X_\Sigma)\ \wedge\ (D^X_\Sigma),
\]
one has $d(\mathcal{M}_0,\mathcal{N})\ge k_X$.
\end{enumerate}
Hence each toggle $\mathcal{M}(X)$ changes the \emph{fewest possible} ground facts (in Hamming distance from $\mathcal{M}_0$)
needed to flip exactly one core axiom while preserving $T_0$ and the remaining cores and witnessing the designated reopened degeneracy.
\end{lemma}

\begin{proof}
We argue case-by-case, using that $T_0$ fixes $\mathrm{Sub}(d_2,d_1)$, $\mathrm{Eval}(d_k,a_i)$ for all $k,i$, and the required redescription-chain facts.

\textbf{$X=\mathbf{I}$.}
In $\mathcal{M}_0$, $\mathrm{About}(e,t_1)$ is false for every evidence $e$.
But $(D^{\mathbf{I}}_\Sigma)$ requires $\mathrm{About}(e',t_1)$ for some evidence $e'$ with a redescription link.
Thus any such $\mathcal{N}$ must differ from $\mathcal{M}_0$ on at least one ground $\mathrm{About}$-tuple, so $d(\mathcal{M}_0,\mathcal{N})\ge 1$.
The model $\mathcal{M}(\mathbf{I})$ differs from $\mathcal{M}_0$ only by adding $\mathrm{About}(e_2,t_1)$, hence achieves distance $1$.

\textbf{$X=\mathbf{C}$.}
Since $(D^{\mathbf{C}}_\Sigma)$ requires $\mathrm{Red}(e_0,e_1)$ and $\mathrm{Red}(e_1,e_2)$ (both fixed true by $T_0$) but $\neg\mathrm{Red}(e_0,e_2)$,
and since $\mathrm{Red}^{\mathcal{M}_0}(e_0,e_2)$ is true, any witness must remove at least that one ground $\mathrm{Red}$-tuple.
Thus $d(\mathcal{M}_0,\mathcal{N})\ge 1$, and $\mathcal{M}(\mathbf{C})$ achieves equality by deleting exactly $(e_0,e_2)$.

\textbf{$X=\mathbf{S}$.}
$(D^{\mathbf{S}}_\Sigma)$ explicitly requires $sv(d_2,a_3)=v_1$ while $sv^{\mathcal{M}_0}(d_2,a_3)=v_0$.
Therefore any witness must change at least one function-value of $sv$, giving $d(\mathcal{M}_0,\mathcal{N})\ge 1$.
The model $\mathcal{M}(\mathbf{S})$ changes exactly the single value $sv(d_2,a_3)$, so the bound is tight.

\textbf{$X=\mathbf{M}$.}
$(D^{\mathbf{M}}_\Sigma)$ requires $sv(d,a_1)=v_1$ for some domain $d$, whereas $\mathcal{M}_0$ has $sv(d,a_1)=v_0$ for all $d$.
Thus at least one $sv$-value must change. But $(S_\Sigma)$ is preserved in this case and, since $\mathrm{Sub}(d_2,d_1)$ and all evaluations hold by $T_0$,
$(S_\Sigma)$ forces $sv(d_1,a_1)=sv(d_2,a_1)$. Hence changing one domain forces changing both: $d(\mathcal{M}_0,\mathcal{N})\ge 2$.
$\mathcal{M}(\mathbf{M})$ changes exactly these two values, so $k_{\mathbf{M}}=2$ is minimal.

\textbf{$X=\mathbf{B}$.}
$(D^{\mathbf{B}}_\Sigma)$ requires $sv(d,a_3)=v_h$ for some domain $d$.
Preserving $(S_\Sigma)$ and using $T_0$ yields $sv(d_1,a_3)=sv(d_2,a_3)$; thus $sv(d,a_3)=v_h$ for one domain forces it for both.
Since $\mathcal{M}_0$ assigns $v_0$ everywhere, at least two $sv$-graph entries must change, so $d(\mathcal{M}_0,\mathcal{N})\ge 2$.
$\mathcal{M}(\mathbf{B})$ achieves this bound.

\textbf{$X=\mathbf{R}$.}
$(D^{\mathbf{R}}_\Sigma)$ requires $\mathrm{Same}(a_2,a_3)$, but $\mathrm{Same}^{\mathcal{M}_0}=\emptyset$, so at least one ground $\mathrm{Same}$-tuple must change.
It also requires $sv(d,a_3)=v_1$ for some domain $d$, while $\mathcal{M}_0$ has $sv(d,a_3)=v_0$ for all $d$.
Preserving $(S_\Sigma)$ again forces the $sv$-change to occur in \emph{both} domains, contributing at least $2$ function-graph changes.
Hence $d(\mathcal{M}_0,\mathcal{N})\ge 1+2=3$, and $\mathcal{M}(\mathbf{R})$ achieves equality.

\textbf{$X=\mathbf{A}$.}
$(D^{\mathbf{A}}_\Sigma)$ requires $\neg\mathrm{LT}(\tau,score(a_0))$ and $\mathrm{LT}(\tau,score(a_3))$.
In $\mathcal{M}_0$, $score(a_0)=score(a_3)=v_0$ and $\neg\mathrm{LT}(\tau,v_0)$ holds by $T_0$.
Thus achieving $\mathrm{LT}(\tau,score(a_3))$ while keeping $\neg\mathrm{LT}(\tau,score(a_0))$ forces $score(a_3)\neq score(a_0)$,
hence at least one $score$-graph entry must change: $d(\mathcal{M}_0,\mathcal{N})\ge 1$.
Moreover, the defining conjunct in $(D^{\mathbf{A}}_\Sigma)$ enforces, for some domain $d$, the threshold equivalence
$sv(d,x)=v_1\leftrightarrow \mathrm{LT}(\tau,score(x))$ for all evaluated acts $x$.
Since $\mathrm{LT}(\tau,score(a_3))$ holds, this forces $sv(d,a_3)=v_1$, whereas $\mathcal{M}_0$ has $sv(d,a_3)=v_0$.
Preserving $(S_\Sigma)$ forces the $sv$-change to occur in both domains, contributing at least $2$ $sv$-graph changes.
Therefore $d(\mathcal{M}_0,\mathcal{N})\ge 1+2=3$.
$\mathcal{M}(\mathbf{A})$ achieves this bound by changing exactly $score(a_3)$ and the two corresponding $sv$-values.
\end{proof}


\begin{remark}[Uniformity guarantee]
The definition of $\Theta$ as a \emph{single-valued} toggle set forces the ``flip exactly one core axiom at a time'' condition:
there is no parameter setting in this appendix in which two distinct perturbations are simultaneously active.
\end{remark}

\end{document}
